\begin{pastexam}[label={px:ncku-114-prog-05}]{成大資訊聯招 114 學年度 第 5 題}
（複選；至少有二個正確答案）

{[4\%]} Please successively insert the data pairs containing the following keys into an empty height-based min leftist tree: 10, 20, 30, 40, 50, 60, 70, 5, 4, 15, 16. After each insertion, the tree should still be a min leftist tree. Which of the following descriptions are correct for the resultant tree?

\textbf{Note: In a tree, each step from top to bottom is called the level of a tree. The level count starts with 1 and increments by 1 at each level or step.}
\begin{enumerate}[label=(\Alph*)]
\item The node containing key=10 has left child with key=30 and right child with key=20.
\item The node containing key=10 and the node containing key=16 are siblings.
\item After repeatedly performing the delete-min operation twice on the leftist tree, the node containing 20 is in the left subtree of the root.
\item The difference of height between the left and right subtrees of the root is 2.
\item The 3rd level of the tree has four nodes.
\end{enumerate}
\end{pastexam}

\begin{pastsoln}
答案：(A)、(D)。

\textbf{規則}（資結書第~8.5~節、定理~8.2）。 高度偏左（height-based）的最小左傾樹：每個節點記 $s(x)$ = 從 $x$ 走到
空子樹的最短距離（空子樹 $s = 0$），要求每個節點 $s(\text{左子}) \ge s(\text{右子})$，於是 $s(x) = s(\text{右子}) + 1$。
插入 $k$ 就是把單一節點的樹和原樹合併（meld）：比較兩個根，較小的當根，它的右子樹和另一棵樹遞迴合併；回來後如果左子的
$s$ 比右子小，就交換左右，再更新 $s$。 下面用 $x(L, R)$ 表示根 $x$、左子樹 $L$、右子樹 $R$，$-$ 是空的。

\textbf{依序插入}：
\begin{center}
\small
\begin{tabular}{rl}
\toprule
插入 & 結果 \\
\midrule
10 & $10$ \\
20 & $10(20, -)$ \\
30 & $10(20, 30)$：$30$ 和右子樹（空）合併放右邊，兩邊 $s$ 都是 $1$，不交換 \\
40 & $10(20, 30(40, -))$：$40$ 併進右子樹 $30$，$30$ 的右邊是 $40$，交換到左邊 \\
50 & $10(30(40, 50), 20)$：$50$ 併進 $30$ 得 $30(40, 50)$，$s = 2$ 大於左子 $20$ 的 $1$，交換 \\
60 & $10(30(40, 50), 20(60, -))$ \\
70 & $10(30(40, 50), 20(60, 70))$ \\
5 & $5(10(\ldots), -)$：$5$ 比 $10$ 小，$5$ 當根，整棵舊樹併進它的右子樹後換到左邊 \\
4 & $4(5(10(\ldots), -), -)$ \\
15 & $4(5(10(\ldots), -), 15)$：$15$ 併進 $4$ 的右子樹（空），$s(5) = 1 \ge s(15) = 1$，不交換 \\
16 & $4(5(10(\ldots), -), 15(16, -))$：$16$ 併進 $15$ 的右子樹後換到左邊 \\
\bottomrule
\end{tabular}
\end{center}
其中 $10(\ldots) = 10(30(40, 50), 20(60, 70))$。 結果：
\begin{center}
\begin{tikzpicture}[x=0.75cm,y=0.9cm,v/.style={circle,draw,inner sep=0pt,minimum size=6mm,font=\small}]
\node[v] (n4) at (5,0) {4};
\node[v] (n5) at (3,-1) {5}; \node[v] (n15) at (7.5,-1) {15};
\node[v] (n10) at (3,-2) {10}; \node[v] (n16) at (7,-2) {16};
\node[v] (n30) at (1.5,-3) {30}; \node[v] (n20) at (4.5,-3) {20};
\node[v] (n40) at (0.8,-4) {40}; \node[v] (n50) at (2.2,-4) {50}; \node[v] (n60) at (3.8,-4) {60}; \node[v] (n70) at (5.2,-4) {70};
\draw (n4)--(n5) (n4)--(n15) (n5)--(n10) (n15)--(n16) (n10)--(n30) (n10)--(n20) (n30)--(n40) (n30)--(n50) (n20)--(n60) (n20)--(n70);
\node[font=\footnotesize,anchor=west] at (8.3,0) {第 1 層};
\node[font=\footnotesize,anchor=west] at (8.3,-1) {第 2 層};
\node[font=\footnotesize,anchor=west] at (8.3,-2) {第 3 層};
\node[font=\footnotesize,anchor=west] at (8.3,-3) {第 4 層};
\node[font=\footnotesize,anchor=west] at (8.3,-4) {第 5 層};
\end{tikzpicture}
\end{center}
（$5$、$16$ 都只有左子節點。） 檢查左傾性質：$s(4) = s(15) + 1 = 2$，$s(5) = 1$、$s(15) = 1$，$s(10) = 3$，……每個節點
左子的 $s$ 都不小於右子。

逐項檢查：
\begin{itemize}
\item (A) 對。 $10$ 的左子是 $30$、右子是 $20$（插入 $50$ 時交換過）。
\item (B) 錯。 $10$ 的父節點是 $5$，$16$ 的父節點是 $15$，不是兄弟。
\item (C) 錯。 刪除最小值就是拿掉根、把左右子樹合併。 第一次：合併 $5(10(\ldots), -)$ 和 $15(16, -)$，$5$ 當根，右子樹變成
$15(16, -)$，得 $5(10(\ldots), 15(16, -))$。 第二次：合併 $10(30(40, 50), 20(60, 70))$ 和 $15(16, -)$，$10$ 當根，右子樹
$20(60, 70)$ 和 $15(16, -)$ 合併得 $15(20(60, 70), 16)$，其 $s = 2$，和左子 $30(40, 50)$ 的 $s = 2$ 相等，不交換。 結果
$10(30(40, 50), 15(20(60, 70), 16))$：$20$ 在根的\textbf{右}子樹裡。
\item (D) 對。 左子樹 $5 \to 10 \to 30 \to 40$ 有 $4$ 層，右子樹 $15 \to 16$ 有 $2$ 層，相差 $2$（用邊數算是 $3$ 和 $1$，差一樣是 $2$）。
\item (E) 錯。 第 $3$ 層只有 $10$ 和 $16$ 兩個節點。
\end{itemize}
\end{pastsoln}
